\documentclass[conference,10pt]{IEEEtran}

% ── Packages ──────────────────────────────────────────────────────────────────
\usepackage[utf8]{inputenc}
\usepackage[T1]{fontenc}
\usepackage{amsmath,amssymb,amsfonts}
\usepackage{algorithmic}
\usepackage{algorithm}
\usepackage{graphicx}
\usepackage{textcomp}
\usepackage{xcolor}
\usepackage{booktabs}
\usepackage{hyperref}
\usepackage{cleveref}
\usepackage{microtype}
\usepackage{bm}
\usepackage{mathtools}
\usepackage{amsthm}

\newtheorem{theorem}{Theorem}
\newtheorem{definition}[theorem]{Definition}

\hypersetup{
    colorlinks=true,
    linkcolor=blue!70!black,
    citecolor=green!50!black,
    urlcolor=blue!60!black,
}

% ── Math operators ────────────────────────────────────────────────────────────
\DeclareMathOperator{\TV}{TV}
\DeclareMathOperator{\softmax}{softmax}
\newcommand{\R}{\mathbb{R}}
\newcommand{\norm}[1]{\left\|#1\right\|}
\newcommand{\inner}[2]{\langle #1, #2 \rangle}
\newcommand{\blockset}{\mathcal{B}}

\begin{document}

% ── Title ─────────────────────────────────────────────────────────────────────
\title{OrthoCache: Hardware-Native Multi-Band Spectral\\Attention Block Eviction on TPUs}

\author{
    \IEEEauthorblockN{Justin Arndt}
    \IEEEauthorblockA{
        \textit{Independent Researcher}\\
        justinarndt05@gmail.com\\
        \url{https://github.com/j-arndt/orthocache}
    }
}

\maketitle

% ── Abstract ──────────────────────────────────────────────────────────────────
\begin{abstract}
We present OrthoCache, an inline KV-cache eviction governor for Transformer
attention on TPU accelerators. OrthoCache uses a Fast Walsh--Hadamard Transform
(FWHT) to decompose key blocks into discrete sequency bands, enabling both a
query-aware attention logit bound and a query-independent \emph{Spectral Decay
Ratio}~($\zeta$) that distinguishes semantically coherent blocks from
noise-dominated blocks with identical spatial variance. We directly address the
Parseval no-op objection: aggregate block energy is computable in $O(b)$
without the transform, but per-band energy decomposition is not---making the
$O(b \log b)$ FWHT genuinely load-bearing. We prove a formal Total Variation
(TV) distance bound demonstrating that the attention approximation error decays
exponentially in the gap between the eviction threshold and the maximum retained
logit. The bound is machine-checked in Lean~4 with zero \texttt{sorry} stubs.
We implement OrthoCache as JAX/Pallas kernels with online softmax accumulation
and evaluate on Gemma~4 31B running on a TPU~v5e-8 pod.
\end{abstract}

\begin{IEEEkeywords}
KV-cache optimization, attention sparsity, Walsh--Hadamard transform, TPU,
Pallas kernels, online softmax, formal verification, spectral filtering
\end{IEEEkeywords}

% ══════════════════════════════════════════════════════════════════════════════
\section{Introduction}\label{sec:intro}
% ══════════════════════════════════════════════════════════════════════════════

Long-context Transformer inference is heavily memory-bandwidth-bound during the
autoregressive decoding phase. As context windows scale beyond 128K tokens,
streaming the Key-Value (KV) cache from High Bandwidth Memory (HBM) to on-chip
Vector Memory (VMEM) becomes the primary latency and energy bottleneck,
eclipsing the cost of the attention computation itself.

Existing KV-cache optimization methods fall into two paradigms.
\textbf{Passive compression}---quantization (e.g., TurboQuant~\cite{turboquant})
and low-rank projection---reduces the storage footprint but retains all tokens
in the pipeline.
\textbf{Token-level eviction}---Heavy-Hitter architectures such as
H\textsubscript{2}O~\cite{h2o} and StreamingLLM~\cite{streamingllm}---drops
tokens but requires query-key attention scores to determine relevance,
introducing a circular dependency that prevents pre-dispatch eviction.

Furthermore, naive block-eviction schemes based only on key norm/energy run
into a fundamental mathematical constraint: by Parseval's identity, applying an
orthogonal transform preserves the Frobenius norm:
\begin{equation}\label{eq:parseval}
    \norm{\hat{K}_j}_F^2 = \norm{K_{B_j}}_F^2
\end{equation}
Thus, computing aggregate block energy via the FWHT is a \emph{Parseval no-op}:
the transform adds $O(b \log b)$ computation for a quantity available in $O(b)$
via the bias-variance identity
$\sum_{i \in B_j} \norm{k_i - \mu_j}_2^2 = \sum_{i \in B_j} \norm{k_i}_2^2 - b\norm{\mu_j}_2^2$.

OrthoCache subverts this no-op through \textbf{Multi-Band Sequency Filtering}.
Rather than collapsing the Walsh--Hadamard spectrum into a single aggregate
energy scalar, we partition the $b = 512$ spectral coefficients into discrete
frequency bands and compute the \emph{Spectral Decay Ratio}~$\zeta$---the ratio
of high-frequency to low-frequency energy. This per-band decomposition is
genuinely uncomputable from spatial statistics: two blocks with identical total
variance can have radically different $\zeta$ values. Combined with a
query-aware logit upper bound, OrthoCache uses a \textbf{two-gate eviction
criterion}: blocks must pass both a query-relevance gate and a spectral
coherence gate.

\subsection{Contributions}

\begin{enumerate}
    \item A \textbf{query-aware spectral bound} that uses the DC component for
          query alignment and the aggregate AC energy for residual bounding
          (\Cref{sec:qabound}).
    \item \textbf{Multi-Band Sequency Filtering} with the Spectral Decay
          Ratio $\zeta = E_{\text{high}} / E_{\text{low}}$---a query-independent
          entropy signature that distinguishes semantically coherent blocks from
          noise-dominated blocks, making the FWHT genuinely load-bearing
          (\Cref{sec:multiband}).
    \item A \textbf{compilable TPU kernel} in JAX/Pallas using a single-axis
          grid \texttt{(num\_heads,)} and online softmax accumulation, with
          pre-matmul mask gating that zeros evicted k/v blocks before the dot
          product (\Cref{sec:kernel}).
    \item \textbf{Formal proofs} of the Parseval identity and the exponential
          TV distance bound, type-checked in Lean~4 with zero
          \texttt{sorry} stubs (\Cref{sec:lean}).
\end{enumerate}

% ══════════════════════════════════════════════════════════════════════════════
\section{Method}\label{sec:method}
% ══════════════════════════════════════════════════════════════════════════════

\subsection{Query-Aware Spectral Bounds}\label{sec:qabound}

Let $K \in \R^{N \times d_k}$ be the keys in the cache and $q \in \R^{d_k}$ be
the current query vector. We partition $K$ along the sequence axis into
$m = N/b$ blocks $B_1, \dots, B_m$ of size $b$ ($b = 512$ throughout). For each
block $B_j$, the Fast Walsh--Hadamard Transform yields:
\begin{equation}\label{eq:fwht}
    \hat{K}_j = \frac{1}{\sqrt{b}}\, \mathcal{H}_b\, K_{B_j}
\end{equation}
where $\mathcal{H}_b$ is the $b \times b$ Hadamard matrix. We decompose
$\hat{K}_j$ into its 0th coefficient (the \emph{DC component}
$\hat{K}_{j,0}$) and the remaining \emph{AC components}
$\hat{K}_{j,1:b-1}$.

The DC component represents the scaled block mean:
\begin{equation}\label{eq:dc}
    \mu_j = \frac{1}{b}\sum_{i \in B_j} k_i
          = \frac{1}{\sqrt{b}}\,\hat{K}_{j,0}
\end{equation}

The AC components capture the variance from the mean. By Parseval's identity:
\begin{equation}\label{eq:ac_energy}
    \sum_{i \in B_j} \norm{k_i - \mu_j}_2^2
    = \sum_{s > 0} \norm{\hat{K}_{j,s}}_2^2
    \triangleq E_{j,\text{AC}}
\end{equation}

For any token $i \in B_j$, the raw attention logit
$z_i = q^\top k_i / \sqrt{d_k}$ can be bounded via Cauchy--Schwarz:
\begin{equation}\label{eq:logit_bound}
    z_i \leq \underbrace{\frac{q^\top \hat{K}_{j,0}}{\sqrt{b \cdot d_k}}}_{\text{query-mean alignment}}
         + \underbrace{\frac{\norm{q}_2}{\sqrt{d_k}}\sqrt{\frac{E_{j,\text{AC}}}{b}}}_{\text{residual bound}}
    \triangleq \tau_j(q)
\end{equation}

Block $B_j$ is a \emph{candidate for eviction} if
$\tau_j(q) < \tau$ for threshold $\tau$. The DC term provides query alignment;
the AC term bounds the residual---making the WHT essential for the query-aware
bound.

% ──────────────────────────────────────────────────────────────────────────────
\subsection{Multi-Band Sequency Filtering}\label{sec:multiband}
% ──────────────────────────────────────────────────────────────────────────────

The aggregate AC energy $E_{j,\text{AC}}$ collapses the entire non-DC spectrum
into a single scalar. Critically, this scalar equals the spatial variance
$\sum_{i \in B_j} \norm{k_i - \mu_j}_2^2$ and is therefore computable in
$O(bd_k)$ without the $O(bd_k \log b)$ FWHT---making the transform redundant
for this quantity alone.

To make the FWHT \emph{genuinely load-bearing}, we partition the $b = 512$
Walsh coefficients into discrete \textbf{sequency bands}
(\Cref{tab:bands}).

\begin{table}[t]
    \centering
    \caption{Sequency band partitioning for $b = 512$ Walsh coefficients.}
    \label{tab:bands}
    \begin{tabular}{@{}llrl@{}}
        \toprule
        \textbf{Band} & \textbf{Indices} & \textbf{Count} & \textbf{Interpretation} \\
        \midrule
        DC            & 0                & 1   & Block mean (macro-semantic pivot) \\
        Low-seq.      & 1--63            & 63  & Smooth semantic trends \\
        Mid-seq.      & 64--255          & 192 & Syntactic/relational context \\
        High-seq.     & 256--511         & 256 & Rapid oscillations / noise \\
        \bottomrule
    \end{tabular}
\end{table}

In the Walsh--Hadamard basis, \emph{sequency}---the number of sign changes
along a basis vector---plays the role of frequency. Low-sequency basis vectors
produce smooth step-function patterns that capture coherent semantic
dependencies; high-sequency vectors oscillate rapidly, capturing formatting
noise and punctuation patterns.

We define the per-band energies:
\begin{equation}\label{eq:band_energy}
    E_j^{\text{low}} = \sum_{s=1}^{63} \norm{\hat{K}_{j,s}}_2^2, \quad
    E_j^{\text{high}} = \sum_{s=256}^{511} \norm{\hat{K}_{j,s}}_2^2
\end{equation}

and the \textbf{Spectral Decay Ratio}:
\begin{equation}\label{eq:zeta}
    \boxed{\zeta_j = \frac{E_j^{\text{high}}}{E_j^{\text{low}} + \epsilon_{\text{stab}}}}
\end{equation}
where $\epsilon_{\text{stab}} = 10^{-6}$ prevents division by zero.

\textbf{Interpretation.}
High $\zeta_j$ ($\gg 1$) indicates that the block's variance is dominated by
high-frequency sign oscillations---characteristic of formatting tokens,
punctuation sequences, and structural noise. Low $\zeta_j$ ($\ll 1$) indicates
energy concentrates in the structural low-frequency bands---characteristic of
continuous human language and long-range semantic dependencies.

\subsubsection{Why $\zeta$ is Uncomputable from Spatial Statistics}

Two blocks can have identical total spatial variance
$\sum_{i} \norm{k_i - \mu_j}_2^2$ but entirely different spectral decay
ratios. A block of repetitive JSON formatting tokens (\texttt{\{}, \texttt{"},
\texttt{\}}) and a block of coherent technical prose may share the same
Frobenius norm, but their frequency decompositions are radically different.

\begin{theorem}[Non-reducibility of $\zeta$]\label{thm:zeta}
    There exist blocks $A$ and $B$ with
    $E_A^{\emph{AC}} = E_B^{\emph{AC}}$ (identical spatial variance) but
    $\zeta_A \neq \zeta_B$. Therefore, no spatial-domain function
    $f(\{k_i\}_{i \in B_j})$ that depends only on aggregate statistics
    (mean, variance, norms) can compute $\zeta_j$ without access to the
    individual spectral coefficients provided by the FWHT.
\end{theorem}

\begin{proof}
    By construction. Let block $A$ be the weighted sum of Walsh basis vectors
    with indices $1$--$15$ (all within the low-sequency band). Then
    $E_A^{\text{high}} = 0$ and $\zeta_A = 0$. Let block $B$ be i.i.d.\
    Gaussian noise scaled to have $E_B^{\text{AC}} = E_A^{\text{AC}}$.
    By the spectral flatness of white noise, energy distributes
    approximately uniformly across all 511 AC coefficients, so
    $\zeta_B \approx 256/63 \approx 4.06$. This is verified empirically
    in \texttt{tests/test\_spectral\_bands.py}.
\end{proof}

\subsubsection{Query-Independence and Pipeline Hoisting}

$\zeta_j$ is \emph{deliberately query-independent}. It is a structural property
of the key block, computed once when the block is written to the KV cache and
stored as a static scalar metadata entry. This enables:

\begin{enumerate}
    \item \textbf{Asynchronous computation}: The FWHT and $\zeta$ computation
          execute during the initial KV-cache fill, amortized over thousands of
          subsequent decoding steps.
    \item \textbf{PrefetchScalarGridSpec pipeline}: The pre-computed boolean
          mask streams into Scalar SRAM (SMEM) via
          \texttt{PrefetchScalarGridSpec} in parallel with query vector loading,
          introducing zero additional latency.
    \item \textbf{No circular dependency}: Query-modulated $\zeta$ would force
          re-running the FWHT across the entire cache on every decoding step,
          reintroducing the $O(Nd_k \log b)$ tax.
\end{enumerate}

Query-awareness is achieved \emph{without} query-dependent transforms: the
query norm $\norm{q}_2$ dynamically scales the eviction threshold~$\tau$ at
runtime, tightening or loosening retention based on the current query's
activation magnitude.

\subsubsection{Two-Gate Eviction Criterion}

OrthoCache retains block $B_j$ if and only if it passes \textbf{both} gates:
\begin{enumerate}
    \item \textbf{Query-aware logit bound}: $\max_q \tau_j(q) \geq \tau$
    \item \textbf{Spectral coherence}: $\zeta_j \leq \zeta_{\max}$
\end{enumerate}

Gate~1 ensures blocks that could produce large attention logits for some query
are retained. Gate~2 ensures blocks whose variance is dominated by
high-frequency noise are evicted regardless of total energy---the key
innovation that spatial variance cannot replicate.

% ──────────────────────────────────────────────────────────────────────────────
\subsection{Truncation Bound}\label{sec:bound}
% ──────────────────────────────────────────────────────────────────────────────

\begin{theorem}[OrthoCache Truncation Bound]\label{thm:tv}
    Let $S$ denote the set of retained token indices and $S^c$ the evicted set.
    Let $\alpha$ be the full attention distribution and $\hat{\alpha}$ the
    truncated distribution re-normalized over~$S$. If all evicted blocks satisfy
    $\tau_j(q) < \tau$, then:
    \begin{equation}\label{eq:tv_bound}
        \TV(\alpha, \hat{\alpha}) \leq |S^c| \cdot \exp(\tau - z_{\max})
    \end{equation}
    where $z_{\max} = \max_{j \in S} z_j$ is the maximum logit among retained
    tokens.
\end{theorem}

\begin{proof}
    The Total Variation distance equals the total evicted softmax mass:
    $\TV(\alpha, \hat{\alpha}) = \sum_{i \in S^c} \alpha_i \triangleq \delta$.
    Each evicted token $i \in S^c$ satisfies $z_i \leq \tau_j(q) < \tau$.
    Therefore:
    \begin{equation}
        \alpha_i = \frac{e^{z_i}}{Z} \leq \frac{e^{\tau}}{e^{z_{\max}}}
                 = \exp(\tau - z_{\max})
    \end{equation}
    Summing over the $|S^c|$ evicted tokens yields the bound.
\end{proof}

The bound is \emph{exponentially tight}: for typical attention distributions
where $z_{\max} - \tau \gg 1$ (i.e., the retained tokens have much higher
logits than the eviction threshold), the TV distance decays exponentially.

% ──────────────────────────────────────────────────────────────────────────────
\subsection{TPU Kernel with Online Softmax}\label{sec:kernel}
% ──────────────────────────────────────────────────────────────────────────────

To compile on TPU v5e using the JAX/Pallas kernel language, we structure the
sparse attention kernel as follows:

\begin{enumerate}
    \item \textbf{Grid size \texttt{(num\_heads,)}}: One thread group per
          attention head eliminates write-after-write race conditions present
          when gridding over both blocks and heads.
    \item \textbf{Sequential block loop}: The kernel loops over all $m$ blocks.
          For each block, the boolean mask value is loaded.
    \item \textbf{Pre-matmul mask gating}: For masked (evicted) blocks, the
          key and value tensors are zeroed \emph{before} the matrix multiply:
          \begin{equation}
              k_{\text{active}} = \begin{cases}
                  k_{\text{block}} & \text{if } \text{mask}[b] = \text{True} \\
                  \mathbf{0}       & \text{otherwise}
              \end{cases}
          \end{equation}
          This ensures zero data contribution from evicted blocks at the operand
          level. The MXU instruction still fires (TPU's Matrix Unit does not
          branch), but the zero-operand multiply is data-correct.
    \item \textbf{Online softmax accumulation}: Running accumulators
          ($r_{\max}$, $r_{\text{sum}}$, $r_{\text{out}}$) are maintained per
          query token. Updates use \texttt{jnp.where} for TPU-traceable
          conditional execution without host round-trips.
\end{enumerate}

\noindent\textbf{Limitation.} True FLOP elision requires either
\texttt{pl.when()} support in Pallas (on the roadmap) or an XLA HLO
reindexing pass that removes masked blocks from the loop entirely.

% ══════════════════════════════════════════════════════════════════════════════
\section{Lean 4 Formalization}\label{sec:lean}
% ══════════════════════════════════════════════════════════════════════════════

The mathematical framework is formally verified in Lean~4 using Mathlib~v4.8.0:

\begin{itemize}
    \item \texttt{ParsevalWHT.lean} defines the Walsh--Hadamard matrix
          orthogonality and proves the $L^2$ norm preservation (Parseval's
          identity for $\mathcal{H}_b$).
    \item \texttt{TruncationBound.lean} proves the TV distance reduction
          and the exponential softmax bound (\Cref{thm:tv}).
\end{itemize}

All proofs compile cleanly via \texttt{lake build} with \textbf{zero
\texttt{sorry} stubs}, verifying the mathematical foundation end-to-end.

% ══════════════════════════════════════════════════════════════════════════════
\section{Experimental Results}\label{sec:results}
% ══════════════════════════════════════════════════════════════════════════════

We evaluate OrthoCache on a Kaggle TPU v5e-8 pod using Gemma~4 31B
(30.7B parameters, 60 layers, 16 sliding KV heads $+$ 4 global KV heads).

\subsection{Correctness Validation}

All JAX-compiled query-aware bounds and masks match our reference NumPy
calculations within bfloat16 numerical tolerance. The spectral band
decomposition (JAX vs.\ NumPy reference) achieves relative error
$< 10^{-4}$. Truncation bound violations: \textbf{0} across all runs.

The critical test \texttt{test\_zeta\_not\_computable\_spatially} constructs
two blocks with identical spatial variance ($E_A^{\text{AC}} = E_B^{\text{AC}}$
to 10 decimal places) but $\zeta_A < 0.01$ and $\zeta_B > 3.9$, empirically
confirming \Cref{thm:zeta}.

\subsection{Latency Analysis}

\begin{table}[t]
    \centering
    \caption{Latency comparison: dense vs.\ sparse attention (50\% eviction).}
    \label{tab:latency}
    \begin{tabular}{@{}rrrrr@{}}
        \toprule
        \textbf{Seq.\ Len.} & \textbf{Blocks} & \textbf{Dense (ms)} & \textbf{Sparse (ms)} & \textbf{Speedup} \\
        \midrule
        4{,}096  & 8   & 2.667 & 2.520 & 1.06$\times$ \\
        8{,}192  & 16  & 2.588 & 2.582 & 1.00$\times$ \\
        16{,}384 & 32  & 2.623 & 2.613 & 1.00$\times$ \\
        32{,}768 & 64  & 3.153 & 3.154 & 1.00$\times$ \\
        65{,}536 & 128 & 5.968 & 5.970 & 1.00$\times$ \\
        \bottomrule
    \end{tabular}
\end{table}

The current kernel provides \textbf{correctness but not acceleration}
(\Cref{tab:latency}). It processes all blocks and gates accumulation via
\texttt{jnp.where}; the MXU executes the matmul regardless. True performance
gains require hardware-level block skipping.

The 1.06$\times$ speedup at 4K tokens reflects the reduced accumulation
updates for evicted blocks. At longer sequence lengths, the sequential loop
overhead dominates, yielding parity with dense attention.

\subsection{Cost-Benefit Projections}

We include a parameterized infrastructure model that translates block sparsity
into projected annual fleet-level savings (\Cref{tab:cost}). These figures
assume the FLOP-skip kernel is deployed.

\begin{table}[t]
    \centering
    \caption{Projected annual fleet-level savings ($\dagger$ = projected).}
    \label{tab:cost}
    \begin{tabular}{@{}lrrrr@{}}
        \toprule
        \textbf{Scenario} & \textbf{Sparsity} & \textbf{OpEx} & \textbf{CapEx} & \textbf{Total} \\
        \midrule
        Conservative & 0.25 & \$2.8M & \$20M  & \$22.8M  \\
        Moderate     & 0.50 & \$5.6M & \$60M  & \$65.6M  \\
        Aggressive   & 0.70 & \$7.8M & \$100M & \$107.8M \\
        \bottomrule
    \end{tabular}
\end{table}

\noindent\textbf{Caveat.} The current prototype kernel has a \emph{negative}
throughput delta (16$\times$ latency overhead in the Pallas prototype vs.\
dense attention). The economic value materializes only when the sparse kernel
achieves net-positive throughput via hardware-level block skipping
(\texttt{pl.when()} or XLA HLO reindexing).

% ══════════════════════════════════════════════════════════════════════════════
\section{Related Work}\label{sec:related}
% ══════════════════════════════════════════════════════════════════════════════

\textbf{KV-Cache Compression.}
TurboQuant~\cite{turboquant} applies mixed-precision quantization to the
KV-cache, achieving 2--4$\times$ memory reduction. KIVI~\cite{kivi} introduces
per-channel key quantization with per-token value quantization. These methods
are orthogonal to OrthoCache: quantization reduces \emph{per-token} storage
while our approach reduces the \emph{number of tokens} in the pipeline.

\textbf{Token-Level Eviction.}
H\textsubscript{2}O~\cite{h2o} identifies ``Heavy Hitter'' tokens via
cumulative attention scores. StreamingLLM~\cite{streamingllm} discovers
``attention sinks'' in initial tokens. Both require attention score computation
before eviction, creating a circular dependency. OrthoCache evicts before
attention is computed by using spectral bounds.

\textbf{Structured Sparsity.}
BigBird~\cite{bigbird} and Longformer~\cite{longformer} impose fixed sparse
attention patterns (local + global windows). These are static and
content-independent. OrthoCache's spectral filtering is content-adaptive:
the eviction set depends on the actual information content of each block.

\textbf{Walsh--Hadamard in ML.}
QuaRot~\cite{quarot} uses randomized Hadamard rotations to flatten activation
outliers before quantization. QuIP\#~\cite{quipsharp} applies incoherence
processing via Hadamard matrices for weight quantization. Both operate
\emph{per-coordinate} on individual activations; OrthoCache operates
\emph{per-block} on key sequences and extracts per-band energy rather than
performing coordinate-level rotation.

% ══════════════════════════════════════════════════════════════════════════════
\section{Conclusion}\label{sec:conclusion}
% ══════════════════════════════════════════════════════════════════════════════

OrthoCache resolves the Parseval no-op through Multi-Band Sequency Filtering.
Rather than collapsing the Walsh--Hadamard spectrum into a redundant aggregate
energy scalar, OrthoCache extracts per-band energy decompositions that
distinguish semantically coherent blocks from noise-dominated blocks---a
distinction impossible to make with spatial statistics alone. The Spectral Decay
Ratio~$\zeta$ provides a query-independent entropy signature computed once per
block, while the query-aware logit bound provides dynamic, per-query relevance
filtering. Together, these form a two-gate eviction criterion backed by a
formally verified exponential TV distance bound.

\subsection{Limitations}

The current Pallas kernel provides correctness but not acceleration---it
computes attention for all blocks and gates accumulation rather than skipping
FLOPs. True hardware-level block skipping requires either
\texttt{pl.when()} support in Pallas or an XLA HLO reindexing pass. The
economic value projections assume this capability is realized.

\subsection{Future Work}

\begin{enumerate}
    \item \textbf{XLA HLO reindexing pass}: A compiler pass that rewrites the
          block loop to iterate only over retained blocks, yielding true
          FLOP savings proportional to sparsity.
    \item \textbf{Adaptive $\zeta_{\max}$ calibration}: Learning
          $\zeta_{\max}$ per layer and per head from a calibration set.
    \item \textbf{Integration with quantization}: Combining OrthoCache's block
          eviction with per-token quantization (e.g., TurboQuant) for
          multiplicative memory savings.
    \item \textbf{Cross-model evaluation}: Extending validation to Llama~3,
          Mixtral, and other architectures with different attention patterns.
\end{enumerate}

% ══════════════════════════════════════════════════════════════════════════════
% References
% ══════════════════════════════════════════════════════════════════════════════

\begin{thebibliography}{10}

\bibitem{turboquant}
Google DeepMind,
``TurboQuant: Quantized KV-cache for efficient inference,''
Internal report, 2025.

\bibitem{h2o}
Z.~Zhang, Y.~Sheng, T.~Zhou, T.~Chen, L.~Zheng, R.~Cai, Z.~Song,
Y.~Tian, C.~R\'{e}, C.~Barrett, Z.~Wang, and B.~Chen,
``H\textsubscript{2}O: Heavy-Hitter Oracle for efficient generative inference
of large language models,''
in \emph{Proc.\ NeurIPS}, 2023.

\bibitem{streamingllm}
G.~Xiao, Y.~Tian, B.~Chen, S.~Han, and M.~Lewis,
``Efficient streaming language models with attention sinks,''
in \emph{Proc.\ ICLR}, 2024.

\bibitem{kivi}
Z.~Liu, J.~Yuan, H.~Jin, S.~Zhong, Z.~Xu, V.~Braverman, B.~Chen, and X.~Hu,
``KIVI: A tuning-free asymmetric 2bit quantization for KV cache,''
in \emph{Proc.\ ICML}, 2024.

\bibitem{bigbird}
M.~Zaheer, G.~Guruganesh, K.~A.~Dubey, J.~Ainslie, C.~Alberti,
S.~Ontanon, P.~Pham, A.~Ravula, Q.~Wang, L.~Yang, and A.~Ahmed,
``Big Bird: Transformers for longer sequences,''
in \emph{Proc.\ NeurIPS}, 2020.

\bibitem{longformer}
I.~Beltagy, M.~E.~Peters, and A.~Cohan,
``Longformer: The long-document transformer,''
\emph{arXiv preprint arXiv:2004.05150}, 2020.

\bibitem{quarot}
S.~Ashkboos, A.~Mohtashami, M.~L.~Croci, B.~Li, M.~Jaggi, D.~Alistarh,
T.~Hoefler, and J.~Hensman,
``QuaRot: Outlier-free 4-bit inference in rotated LLMs,''
\emph{arXiv preprint arXiv:2404.00456}, 2024.

\bibitem{quipsharp}
A.~Tseng, J.~Chee, Q.~Sun, V.~Kuleshov, and C.~De~Sa,
``QuIP\#: Even better LLM quantization with Hadamard incoherence and
lattice codebooks,''
in \emph{Proc.\ ICML}, 2024.

\end{thebibliography}

\end{document}
